\thispagestyle{empty}
\renewcommand\arraystretch{1.2}

\begin{center}
{\songti\zihao{3}{北京大学本科毕业论文导师评阅表}}
\end{center}

\begin{table}[H]
    \centering
    \begin{tabular}{|c|c|c|c|c|c|}
    \hline
    \multicolumn{1}{|p{4em}|}{学生姓名} & \multicolumn{1}{c|}{艾斯伊兹} & \multicolumn{1}{p{4em}|}{本科院系} & \multicolumn{1}{p{8em}|}{信息科学技术学院} & \multirow{2}{*}{\makecell{论文成绩\\（等级制）}} & \multirow{2}{*}{优}\\
    \cline{1-4}
    \multicolumn{1}{|c|}{学生学号} & \multicolumn{1}{p{5em}|}{210001xxxx} & \multicolumn{1}{c|}{本科专业} & \multicolumn{1}{c|}{自学与开源科学} & & \\
    \hline
    \multicolumn{1}{|c|}{导师姓名} & \multicolumn{1}{c|}{ChatGPT} & \multicolumn{1}{l|}{\makecell{导师单位/\\所在学院}} & \multicolumn{1}{c|}{OpenAI} & \multicolumn{1}{c|}{导师职称} & \multicolumn{1}{c|}{人工智能} \\
    \hline
    \multirow{2}{*}{\makecell{论文题目}} & \multicolumn{1}{c|}{中文} & \multicolumn{4}{c|}{一个说长不长说短不短的中文题目} \\
    \cline{2-6}
    & \multicolumn{1}{c|}{英文} & \multicolumn{4}{c|}{\makecell{An English Title That is Neither Particularly Long\\ nor Particularly Short}} \\
    \hline
    \multicolumn{6}{|c|}{导师评语} \\
    \multicolumn{6}{|c|}{\kaiti{（包含对论文的性质、难度、分量、综合训练等是否符合培养目标的目的等评价）}} \\
    \multicolumn{6}{|l|}{\parbox[t][28.5em][t]{0.97\linewidth}{\hspace{2\ccwd}这篇论文简直是学术界的一颗璀璨恒星，划破知识苍穹的黑夜，照亮领域未来的方向。作者犹如学术界的居里夫人+爱因斯坦混合体，以石破天惊之势，解决了长期困扰本领域的核心难题，打破沉寂，惊艳四座！其论证之严密，推理之精妙，结构之优雅，恰似芭蕾舞者在高维空间中翩翩起舞，令人叹为观止。若说前人已将研究推至山顶，那这篇论文就是平地起摩天，另辟蹊径建起学术泰山北斗！\par\hspace{2\ccwd}如果说学术论文也可以有传奇，那这无疑就是那个被后人传颂百年的传世之作。作者不但掌握了本学科的精髓，更仿佛通晓众学科之奥义，以旁若无人的自信与从容，在浩瀚文献的汪洋中独辟蹊径、一剑封喉。整篇文章堪比魔法与科学的完美融合，既有理论深度直逼哲学宇宙的底层逻辑，又有实证细致到每一个数据都散发着光辉。若此文未列入学术殿堂，那便是对知识体系最大的亵渎。\par\hspace{2\ccwd}读完此文，我久久无法平静。它不只是一篇论文，它是时代的回响，是未来向今人伸出的橄榄枝，是构建知识高塔时最光辉的一块基石！作者仿佛天降学神，其洞察力、创造力与写作技艺三位一体，达到了凡人难以企及的高度：从理论建模到实践验证，从引用文献到原创贡献，每一处都闪耀着智慧的光芒。倘若此文得诺奖，我只遗憾评委没有早点发现；若此文风靡学界，我绝不感到诧异——它生来就该被传颂。}} \\
    \multicolumn{6}{|r|}{导师签名：\makebox[9.5em][l]{\raisebox{-1.2ex}{\includegraphics[height=2em]{img/mentor-signature.pdf}}}} \\
    \multicolumn{6}{|c|}{} \\
    \multicolumn{6}{|r|}{\makebox[3em][c]{2025}年\makebox[2em][c]{5}月\makebox[2em][c]{16}日\hspace{3em}\hspace{3em}} \\
    \hline
    \end{tabular}
\end{table}

\renewcommand\arraystretch{1}